\section{Result 2: unequal use of the animals did no harm}
\label{sec:zipf}

You pointed out that in real language words are not used equally often, and that this is part
of what a language is. So the first round was run again with unequal use: the most
common animal is mentioned twice as often as the second, three times as often as the third,
and so on down the list. The 16 original animals are the most common ones. With 1,024 animals
they get 45\% of all mentions, and the rarest animal is mentioned 7.5 times less often than it
would be with equal use. These are again 36 runs. To fit the night I cut the time limit from 30
minutes to 5, roughly 270 sets. That turned out not to matter: every run was solved inside it,
the slowest after 256 sets.

\textbf{All 36 runs ended as solved.} The lowest score was 99.1. So the rule is again not met
for any width, and for the same reason: no model ended up worse with more words.

\textbf{The large languages were learned in fewer sets than with equal use.} On 1,024 animals
width 32 needed 45 sets in all three repeats, where it had needed about 128 (table
\ref{tab:effort1}). Width 8 needed 181, 181 and 256 sets, where it had needed 256, 512 and
1,639. The rare
animals were learned too. When width 32 was solved on 1,024 animals, each fact of the rarest
animal had come up only about 90 times. With equal use, each fact had come up more than 1,300
times before the same model was solved.

I had expected the opposite, that rare words would hold a model back. My explanation, which
these runs do not test: the common animals get a model past its early stand-still sooner, and
after that it needs far fewer showings to add another animal.
